\section{Introduction to \LaTeX}

\LaTeX, I had never heard about this term before doing this project,
but when I came to know about it's features, found it excellent. 
\LaTeX{} (pronounced /ˈleɪtɛk/, /ˈleɪtɛx/, /ˈlɑːtɛx/, or /ˈlɑːtɛk/) is a 
document markup language and document preparation system for the \TeX{} 
typesetting  program. Within the typesetting system, its name is styled 
as \LaTeX.

\begin{figure}[h]
	\centering \includegraphics[scale=0.2]{images/latex.png}
	\caption{Logo of \LaTeX}
\end{figure}


Within the typesetting system, its name is styled as \LaTeX. The term 
\LaTeX{} refers only to the language in which documents are written, 
not to the editor used to write those documents. In order to create a 
document in \LaTeX, a .tex file must be created using some form of text 
editor. While most text editors can be used to create a \LaTeX{} document, 
a number of editors have been created specifically for working with \LaTeX.

\LaTeX{} is most widely used by mathematicians, scientists, 
engineers, philosophers, linguists, economists and other scholars in 
academia. As a primary or intermediate format, e.g., translating DocBook 
and other XML-based formats to PDF, \LaTeX{} is used because of the 
high quality of typesetting achievable by \TeX. The typesetting system 
offers programmable desktop publishing features and extensive facilities 
for automating most aspects of typesetting and desktop publishing, 
including numbering and cross-referencing, tables and figures, 
page layout and bibliographies.

\LaTeX{} is intended to provide a high-level language that
accesses the power of \TeX. \LaTeX{} essentially comprises a
collection of \TeX{} macros and a program to process \LaTeX documents. 
Because the \TeX{} formatting commands are very low-level, it is usually 
much simpler for end-users to use \LaTeX{}.


\subsection{Typesetting}
\LaTeX{} is based on the idea that authors should be able to focus on 
the content of what they are writing without being distracted by its 
visual presentation. In preparing a \LaTeX{} document, the author 
specifies the logical structure using familiar concepts such as 
chapter, section, table, figure, etc., and lets the \LaTeX{} system 
worry about the presentation of these structures. It therefore 
encourages the separation of layout from content while still allowing 
manual typesetting adjustments where needed. 


